\section{Etapa 10 --- Ruptura arquitect\'onica y portaci\'on a la web}
\label{sec:et10}
\periodo{21 de julio --- 1\textsuperscript{o} de agosto de 2026 \quad$\cdot$\quad 7 d\'ias, \textasciitilde{}40 \emph{commits}}

Etapa de infraestructura. No agrega capacidad de medici\'on, pero cambia c\'omo
se sostiene todo lo dem\'as ---y contiene un pivote de herramienta importante.

\begin{figure}[ht]
\centering
\includegraphics[width=\textwidth]{flujo_adquisicion_masw.png}
\caption{Flujo funcional que gui\'o la separaci\'on entre firmware, datos y aplicaci\'on
web dentro del superproyecto.}
\end{figure}

\subsection{Del monolito al superproyecto (21 de julio)}

El repositorio \'unico se parti\'o. El historial monol\'itico se preserv\'o en
\texttt{Tesis-legacy} y el repositorio \texttt{Tesis} se reinici\'o con un
\textbf{\emph{commit} ra\'iz sin padres} que fija los subm\'odulos como
repositorios independientes.

\begin{center}\small
\begin{tabularx}{\textwidth}{@{}l X@{}}
\toprule
\textbf{Motivo} & \textbf{Consecuencia} \\
\midrule
El firmware, el modelado y los datos ten\'ian ciclos de vida distintos & Cada uno versiona a su ritmo \\
\addlinespace[2pt]
Los datos crudos (\textasciitilde{}29~GiB) envenenaban el historial & LFS con almacenamiento por carpeta y deduplicaci\'on SHA-256, con manifiesto de ruta/tama\~no/\emph{hash} por archivo \\
\addlinespace[2pt]
Un clon completo era impracticable & Se puede clonar solo lo que se necesita \\
\bottomrule
\end{tabularx}
\end{center}

\subsection{Reorganizaci\'on por prop\'osito (24 de julio)}

Tres d\'ias despu\'es, segunda pasada: \texttt{src/} deja de agrupar por lenguaje
y pasa a agrupar \textbf{por prop\'osito}
---\texttt{firmware/}, \texttt{interfaces/}, \texttt{calculos\_modelados/}---
mediante \texttt{subtree split}, con dos repositorios nuevos y dos renombrados.
\texttt{investigacion/} baj\'o de la ra\'iz a \texttt{docs/investigacion}.

\seleccionado{Agrupar por prop\'osito y no por lenguaje es lo correcto ac\'a
porque \emph{el mismo prop\'osito se implementa en varios lenguajes}: hay
interfaces en Python y en MATLAB, y modelado en Python y en MATLAB. Agrupar
por lenguaje separaba cosas que se usan juntas.}

\pendiente{Dos trampas conocidas de esta arquitectura, que ya costaron trabajo:
\textbf{(1)} un subm\'odulo en \emph{detached HEAD} acepta un
\texttt{push origin main} \emph{sin subir nada}, informando \'exito ---hay que
auditar el \emph{gitlink} contra \texttt{ls-remote}; \textbf{(2)}
\texttt{git submodule deinit -f} vac\'ia el \emph{worktree} y destruye lo no
confirmado. Para mover un subm\'odulo hay que editar solo el \texttt{path} en
\texttt{.gitmodules}.}

\subsection{El pivote de interfaz: de PyQt a web}

Hasta julio la revisi\'on de campo viv\'ia en \texttt{field\_review\_app}, una
aplicaci\'on de escritorio PyQt. Entre el 25 de julio y el 29 se port\'o
entera a la web.

\begin{center}\small
\begin{tabularx}{\textwidth}{@{}l X@{}}
\toprule
25 jul & Arranca el \texttt{PORT\_PLAN}: el servidor Python migra de un \texttt{http.server} artesanal a \textbf{FastAPI} con est\'aticos; pesta\~nas de navegaci\'on y alternancia de tema; visor de se\~nal con \textbf{decimado min/max}. En el camino se corrige un bug del propio \emph{gate} automatizado, que daba falsos vencimientos porque un \emph{pipe} de subproceso no ten\'ia lector \\
\addlinespace[2pt]
26 jul & \textbf{Se congela un \emph{snapshot} de referencia de la app PyQt completa} en \texttt{docs/legacy/pyqt\_field\_review/}, con mapa de l\'ineas de la ventana ``Capturas'', \emph{antes} de empezar a portarla \\
\addlinespace[2pt]
29 jul & Portaci\'on masiva hasta alcanzar paridad funcional \\
\bottomrule
\end{tabularx}
\end{center}

\seleccionado{\textbf{Congelar el original antes de portar} es la decisi\'on que
hizo viable el porte: da un objetivo de paridad verificable en vez de una
reescritura ``a ojo''. El decimado min/max, por su parte, es la t\'ecnica
correcta para dibujar millones de muestras en un \emph{canvas} sin perder los
picos ---que en una traza s\'ismica son exactamente lo que importa.}

\subsubsection{Lo que trajo la portaci\'on (29 de julio)}

\begin{itemize}[itemsep=0.2em]
\item \textbf{Cuarentena reversible} (\texttt{deletion.py}) en vez de borrado
f\'isico: un registro rechazado se puede recuperar. Dado que un descarte
err\'oneo de datos de campo es irreparable, es la elecci\'on correcta.
\item \textbf{Cola \texttt{AnalysisJobs}} persistente para los an\'alisis
pesados, que ya no bloquean la interfaz.
\item Exportaciones reproducibles.
\item \textbf{\texttt{smoke\_test.py} con 58 verificaciones} como condici\'on de
avance: el porte no avanzaba si el \emph{gate} no pasaba.
\end{itemize}

\subsection{Otras l\'ineas del per\'iodo}

\textbf{Conectividad de campo.} Arranca el plan del modo ENLACE del maestro y
un \emph{sketch} descartable (\texttt{MasterFoo}) que valida el uso del
\emph{hotspot} del celular como red de campo.

\descartado{\texttt{MasterFoo} era deliberadamente descartable: se escribi\'o
para responder una pregunta (¿funciona el \emph{hotspot}?), se la respondi\'o y
no se lo integr\'o. Registrarlo evita que alguien lo confunda con c\'odigo de
producci\'on.}

\textbf{Martillo de leva} (1\textsuperscript{o} de agosto). Arranca el dise\~no
de una fuente de impacto repetible ---clave para el arreglo sintetizado de la
Etapa~7, donde la repetibilidad entre disparos \emph{es} la
calidad del dato. Se model\'o en dos herramientas independientes: Simscape
Multibody (MATLAB) y un modelo 2D en Python.

\dato{\textbf{Los dos modelos detectaron errores propios distintos}, lo cual es
justamente el argumento para hacerlos por duplicado: el de Simscape encontr\'o
una \emph{interferencia axial entre la leva y el brazo}, y el de Python una
\emph{violaci\'on de conservaci\'on de energ\'ia} en el criterio de suelta.}

\textbf{PCBs.} Entra el proyecto de placa (KiCad, y luego JitX) como
repositorio propio.

\subsection{Balance de la etapa}

\begin{center}\small
\begin{tabularx}{\textwidth}{@{}l X@{}}
\toprule
\textbf{Se qued\'o} & Superproyecto con subm\'odulos por prop\'osito; LFS con dedup SHA-256 y manifiesto; servidor FastAPI; interfaz de revisi\'on en web; decimado min/max; cuarentena reversible; cola \texttt{AnalysisJobs}; \emph{gate} de 58 verificaciones; martillo de leva modelado por duplicado \\
\addlinespace[3pt]
\textbf{Se descart\'o} & Repositorio monol\'itico; organizaci\'on de \texttt{src/} por lenguaje; \texttt{http.server} artesanal; app de escritorio PyQt (congelada como referencia, no borrada); \emph{sketch} \texttt{MasterFoo}; borrado f\'isico de registros \\
\bottomrule
\end{tabularx}
\end{center}
